\documentclass{article}
\usepackage{amsmath, amssymb, amsthm}
\usepackage{hyperref}
\usepackage{geometry}
\geometry{margin=1in}

\title{Erd\H{o}s Problem \#879}
\date{Last updated: 2025-08-31}
\author{Source: erdosproblems.com}

\begin{document}
\maketitle

\noindent\textbf{Status:} OPEN \quad \textbf{No prize}

\noindent\textbf{Tags:} number theory

\medskip
\noindent\textbf{Problem Statement:}

Call a set $S\subseteq \{1,\ldots,n\}$ admissible if $(a,b)=1$ for all $a\neq b\in S$. Let\[G(n) = \max_{S\subseteq \{1,\ldots,n\}} \sum_{a\in S}a\]and\[H(n)=\sum_{p<n}p+ n\pi(n^{1/2}).\]Is it true that\[G(n) >H(n)-n^{1+o(1)}?\]Is it true that, for every $k\geq 2$, if $n$ is sufficiently large then the admissible set which maximises $G(n)$ contains at least one integer with at least $k$ prime factors?

\bigskip
\noindent\textbf{Remarks:}

Erd\H{o}s and Van Lint proved that\[H(n)-n^{3/2-o(1)}<G(n)<H(n)\]and\[\frac{H(n)-G(n)}{n}\to \infty.\]They proved that $G(n)>H(n)-n^{1+o(1)}$ assuming 'plausible (but hopeless) assumptions about the distribution of primes'. They also prove the second claim when $k=2$.

See also [878].

\bigskip
\noindent\small{Source: \url{https://www.erdosproblems.com/879}}
\end{document}
